\section{Định lượng acetic acid trong giấm}

\subsection{Dùng chuẩn độ để xác định hàm lượng acid}

Kĩ thuật chuẩn độ acid--base đã được học ở Chuyên đề 3. Trong chuyên đề này, chuẩn độ được dùng như một công cụ để xác định lượng acetic acid trong mẫu giấm.

Phản ứng:
\[
\ce{CH3COOH + NaOH -> CH3COONa + H2O}.
\]

Tại điểm tương đương:
\[
n_{\ce{CH3COOH}}=n_{\ce{NaOH}}.
\]

\begin{vidu}
Lấy \SI{10.00}{\milli\liter} dung dịch giấm, pha loãng thành \SI{100.0}{\milli\liter}. Lấy \SI{10.00}{\milli\liter} dung dịch sau pha loãng để chuẩn độ bằng dung dịch \ce{NaOH} \SI{0.100}{\mole\per\liter}. Ba lần chuẩn độ dùng lần lượt $8.10$, $8.00$ và $8.05$ mL \ce{NaOH}. Tính nồng độ mol của acetic acid trong mẫu giấm ban đầu.
\end{vidu}
\begin{loigiai}
Thể tích \ce{NaOH} trung bình:
\[
\overline V=\frac{8.10+8.00+8.05}{3}=\SI{8.05}{\milli\liter}.
\]
\[
n_{\ce{NaOH}}=0.100\times8.05\times10^{-3}=\SI{8.05e-4}{\mole}.
\]

\bangmol{5}
{\ce{CH3COOH} & + & \ce{NaOH} & \rightarrow & \ce{CH3COONa + H2O}}
{8.05\times10^{-4} & & 8.05\times10^{-4} & &}

Số mol acetic acid trong toàn bộ \SI{100.0}{\milli\liter} dung dịch sau pha loãng:
\[
n=8.05\times10^{-4}\times\frac{100.0}{10.00}=\SI{8.05e-3}{\mole}.
\]
Đây cũng là số mol acetic acid có trong \SI{10.00}{\milli\liter} giấm ban đầu, nên
\[
C_M=\frac{8.05\times10^{-3}}{10.00\times10^{-3}}=\SI{0.805}{\mole\per\liter}.
\]

\textbf{Vậy,} nồng độ mol của acetic acid trong mẫu giấm là \SI{0.805}{\mole\per\liter}.
\end{loigiai}

\subsection{Đọc và đánh giá số liệu chuẩn độ}

Khi một phép chuẩn độ được lặp lại nhiều lần, các giá trị gần nhau cho thấy kết quả ổn định hơn. Nếu một giá trị lệch rõ rệt so với các lần còn lại, cần xem xét thao tác thí nghiệm trước khi dùng để tính trung bình.

\begin{ghinho}{Một số nguồn sai số thường gặp}
\begin{itemize}
\item burette còn nước làm dung dịch chuẩn bị pha loãng;
\item đọc thể tích không ngang tầm mắt;
\item nhỏ dung dịch chuẩn quá nhanh gần điểm kết thúc;
\item không trộn đều dung dịch trong bình nón;
\item lấy sai thể tích mẫu hoặc pha loãng không đúng vạch.
\end{itemize}
\end{ghinho}

\begin{vidu}
Trước khi nạp dung dịch \ce{NaOH}, một học sinh chỉ tráng burette bằng nước cất mà không tráng lại bằng dung dịch \ce{NaOH}. Giả sử các thao tác khác đúng. Kết quả tính nồng độ acetic acid trong mẫu giấm sẽ có xu hướng lớn hơn hay nhỏ hơn giá trị thật?
\end{vidu}
\begin{loigiai}
Nước còn trong burette làm dung dịch \ce{NaOH} bị pha loãng. Vì vậy phải dùng thể tích dung dịch lớn hơn để trung hòa cùng một lượng acetic acid.

Nếu vẫn tính bằng nồng độ \ce{NaOH} ghi trên nhãn, số mol \ce{NaOH} được tính sẽ lớn hơn số mol thực tế đã phản ứng, từ đó làm nồng độ acetic acid tính được lớn hơn giá trị thật.

\textbf{Vậy,} kết quả có xu hướng cao hơn giá trị thật.
\end{loigiai}

\begin{tuluyen}
Một mẫu giấm được pha loãng $10$ lần. Lấy \SI{10.00}{\milli\liter} dung dịch sau pha loãng cần trung bình \SI{7.50}{\milli\liter} \ce{NaOH} \SI{0.100}{\mole\per\liter} để trung hòa. Tính nồng độ mol acetic acid trong mẫu giấm ban đầu.
\dapso{\SI{0.750}{\mole\per\liter}.}
\end{tuluyen}
